%! AP Statistics — Unit 2, Lesson 2.1 — Slides (Beamer)
\documentclass[aspectratio=169, 11pt]{beamer}
\usepackage{apstats-beamer}

%=============================================================================
\begin{document}

% ─────────────────────────────────────────────────────────────────────────────
% SLIDE 1 — LESSON TITLE
% ─────────────────────────────────────────────────────────────────────────────
{
\setbeamercolor{background canvas}{bg=navy}
\begin{frame}[plain]
  \begin{tikzpicture}[remember picture, overlay]
    \fill[navylight] (current page.north west)
      rectangle ([yshift=-1.35cm]current page.north east);
  \end{tikzpicture}
  \vspace{2.8cm}
  \hspace{0.55cm}%
  \begin{minipage}{0.9\textwidth}
    {\color{goldacc}\bfseries\small LESSON 2.1}\\[0.5cm]
    {\color{white}\bfseries\LARGE Introducing Statistics:\\[0.25cm]
      Are Variables Related?}\\[1.4cm]
    {\color{skymid}\small AP Statistics~$\cdot$~Unit 2: Exploring Two-Variable Data}
  \end{minipage}
\end{frame}
}

% ─────────────────────────────────────────────────────────────────────────────
% SLIDE 2 — PRIMARY OBJECTIVE
% ─────────────────────────────────────────────────────────────────────────────
{
\setbeamercolor{background canvas}{bg=sky}
\begin{frame}[t, plain]
  \begin{tikzpicture}[remember picture, overlay]
    \fill[navy] (current page.north west)
      rectangle ([yshift=-0.25cm]current page.north east);
  \end{tikzpicture}
  \vspace{0.45cm}\par
  \hspace{0.55cm}%
  \begin{minipage}{0.9\textwidth}

    \sectionlabel{Primary Objective}

    \normalsize
    Students will identify statistical questions about possible relationships
    between two variables and recognize that \textbf{apparent patterns in data
    may be random} --- not every association is real, and even a real association
    does not prove causation (Big Idea: VAR, Skill 1.A).

    \vspace{0.6cm}

    \normalsize\textbf{\color{navy}By the end of this lesson, I will be able to:}

    \smallskip
    \begin{itemize}
      \setlength{\itemsep}{0.3em}\setlength{\topsep}{0.2em}
      \item Identify a statistical question about a possible relationship between
        two variables.
      \item Explain why an apparent association may be random, not real.
      \item Explain why association does not imply causation, and give an example
        of a confounding variable.
    \end{itemize}

  \end{minipage}
\end{frame}
}

% ─────────────────────────────────────────────────────────────────────────────
% SLIDE 3 — HOOK
% ─────────────────────────────────────────────────────────────────────────────
{
\setbeamercolor{background canvas}{bg=hookbg}
\setbeamercolor{itemize item}{fg=white}
\begin{frame}[t, plain]
  \vspace{0.45cm}\par
  \hspace{0.55cm}%
  \begin{minipage}{0.9\textwidth}

    {\color{white}\bfseries\footnotesize\MakeUppercase{Hook}}\par\medskip

    {\color{white}\large\bfseries Two Headlines:}\\[0.3cm]
    {\color{skymid}\itshape
      ``Students who eat breakfast score higher on tests.''}\\[0.2cm]
    {\color{skymid}\itshape
      ``Shoe size and reading ability are strongly associated in children.''}

    \vspace{0.5cm}
    {\color{white}\bfseries Discussion questions:}\par\smallskip
    \begin{itemize}
      \setlength{\itemsep}{0.4em}\setlength{\topsep}{0.1em}
      \item {\color{white}Is the breakfast--grades association real? Random? How would you know?}
      \item {\color{white}Shoe size does not teach children to read.
        What actually explains this association?}
      \item {\color{white}These are the questions Unit 2 is designed to answer.}
    \end{itemize}

  \end{minipage}
\end{frame}
}

% ─────────────────────────────────────────────────────────────────────────────
% SLIDE 4 — FROM ONE VARIABLE TO TWO
% ─────────────────────────────────────────────────────────────────────────────
{
\setbeamercolor{background canvas}{bg=white}
\begin{frame}[t, plain]
  \navyheader{From One-Variable to Two-Variable Questions}
  \hspace{0.4cm}%
  \begin{minipage}{0.93\textwidth}
    \vspace{0.3cm}
    \begin{columns}[T]

      \begin{column}{0.48\textwidth}
        \small
        \textbf{\color{navy}Unit 1 --- One variable at a time}
        \begin{itemize}
          \setlength{\itemsep}{0.35em}
          \item What is the typical GPA of students at our school?
          \item How many hours of sleep do students get per night?
          \item What is the distribution of commute times?
        \end{itemize}
      \end{column}

      \begin{column}{0.48\textwidth}
        \small
        \textbf{\color{navy}Unit 2 --- Two variables together}
        \begin{itemize}
          \setlength{\itemsep}{0.35em}
          \item Is there a relationship between sleep and GPA?
          \item Do athletes sleep more or less than non-athletes?
          \item Does commute time affect job satisfaction?
        \end{itemize}
      \end{column}

    \end{columns}
    \vspace{0.6cm}
    \small\textbf{The new question:}\enspace
    Does knowing one variable tell us something about the other?
  \end{minipage}
\end{frame}
}

% ─────────────────────────────────────────────────────────────────────────────
% SLIDE 5 — APPARENT ASSOCIATIONS MAY BE RANDOM
% ─────────────────────────────────────────────────────────────────────────────
{
\setbeamercolor{background canvas}{bg=greenbg}
\begin{frame}[t, plain]
  \begin{tikzpicture}[remember picture, overlay]
    \fill[greenacc] (current page.north west)
      rectangle ([yshift=-1.35cm]current page.north east);
    \node[anchor=west, text=white, font=\bfseries\large,
          inner xsep=0.55cm]
      at ([yshift=-0.68cm]current page.north west)
      {Key Idea: Apparent Patterns May Be Random (VAR-1.D.1)};
  \end{tikzpicture}
  \vspace{1.05cm}\par
  \hspace{0.4cm}%
  \begin{minipage}{0.93\textwidth}
    \vspace{1em}
    \small
    \begin{columns}[T]

      \begin{column}{0.48\textwidth}
        \textbf{\color{greenacc}Example}\\[4pt]
        In a class of 10 students, the 3 who drink coffee
        all have higher GPAs.\\[8pt]
        Does this prove coffee helps GPA?\\[4pt]
        \textbf{No.} With only 10 students, this pattern
        could easily be random. A larger sample
        might show no relationship at all.
      \end{column}

      \begin{column}{0.48\textwidth}
        \textbf{\color{greenacc}What makes an association more credible?}
        \begin{itemize}
          \setlength{\itemsep}{0.3em}
          \item Larger sample size
          \item Stronger association
          \item Consistent across different groups and settings
        \end{itemize}
        \vspace{8pt}
        \textbf{Bottom line:} Statistics gives us the tools to decide
        whether a pattern is real. We build those tools in Units 5--9.
      \end{column}

    \end{columns}
  \end{minipage}
\end{frame}
}

% ─────────────────────────────────────────────────────────────────────────────
% SLIDE 6 — ASSOCIATION vs. CAUSATION
% ─────────────────────────────────────────────────────────────────────────────
{
\setbeamercolor{background canvas}{bg=white}
\begin{frame}[t, plain]
  \navyheader{Association vs.\ Causation}
  \hspace{0.4cm}%
  \begin{minipage}{0.93\textwidth}
    \vspace{0.3cm}
    \begin{columns}[T]

      \begin{column}{0.48\textwidth}
        \small
        \textbf{\color{navy}Three explanations for an association:}
        \begin{enumerate}
          \setlength{\itemsep}{0.4em}
          \item $X$ causes $Y$
          \item $Y$ causes $X$ (reverse causation)
          \item A third variable causes both $X$ and $Y$
                (\textbf{confounding variable})
        \end{enumerate}
        \vspace{8pt}
        \textbf{\color{navy}Classic example --- fire trucks:}\\
        \textit{Cities that send more trucks have more property damage.}\\[4pt]
        Do trucks cause damage?\\
        \textbf{No.} Larger fires cause both.
        Fire size is the confounding variable.
      \end{column}

      \begin{column}{0.48\textwidth}
        \small
        \textbf{\color{navy}Shoe size \& reading ability}\\[4pt]
        Both increase as children \textbf{grow older}.\\
        Age is the confounding variable.\\[8pt]
        \textbf{\color{navy}Ice cream \& drowning}\\[4pt]
        Both increase in \textbf{summer}.\\
        Season / temperature is the confounding variable.\\[8pt]
        \textit{The association is real in both cases ---
        but causation is not.}
      \end{column}

    \end{columns}
  \end{minipage}
\end{frame}
}

% ─────────────────────────────────────────────────────────────────────────────
% SLIDE 7 — VOCABULARY
% ─────────────────────────────────────────────────────────────────────────────
{
\setbeamercolor{background canvas}{bg=greenbg}
\begin{frame}[t, plain]
  \begin{tikzpicture}[remember picture, overlay]
    \fill[greenacc] (current page.north west)
      rectangle ([yshift=-1.35cm]current page.north east);
    \node[anchor=west, text=white, font=\bfseries\large,
          inner xsep=0.55cm]
      at ([yshift=-0.68cm]current page.north west)
      {Vocabulary \& Key Concepts};
  \end{tikzpicture}
  \vspace{1.05cm}\par
  \hspace{0.4cm}%
  \begin{minipage}{0.93\textwidth}
    \vspace{1em}
    \footnotesize
    \begin{multicols}{2}
      \textbf{\color{greenacc}Association} --- a relationship between two variables
        such that knowing one gives information about the other\\[4pt]
      \textbf{\color{greenacc}Apparent association} --- a pattern in a sample that may
        or may not reflect a real relationship in the population\\[4pt]
      \textbf{\color{greenacc}Causation} --- one variable directly produces a change in
        another; requires a controlled experiment to establish\\[4pt]

      \columnbreak

      \textbf{\color{greenacc}Confounding variable} --- a third variable related to both
        the explanatory and response variable, making an apparent association misleading\\[4pt]
      \textbf{\color{greenacc}Explanatory variable} --- the variable used to explain or
        predict the response; placed on the $x$-axis\\[4pt]
      \textbf{\color{greenacc}Response variable} --- the outcome of interest;
        placed on the $y$-axis
    \end{multicols}
  \end{minipage}
\end{frame}
}

% ─────────────────────────────────────────────────────────────────────────────
% SLIDE 8 — GROUP ACTIVITY
% ─────────────────────────────────────────────────────────────────────────────
{
\setbeamercolor{background canvas}{bg=redbg}
\setbeamercolor{enumerate item}{fg=redacc}
\begin{frame}[t, plain]
  \begin{tikzpicture}[remember picture, overlay]
    \fill[redacc] (current page.north west)
      rectangle ([yshift=-1.35cm]current page.north east);
    \node[anchor=west, text=white, font=\bfseries\large,
          inner xsep=0.55cm, inner ysep=0pt]
      at ([yshift=-0.675cm]current page.north west)
      {Group Activity};
  \end{tikzpicture}
  \vspace{1.05cm}\par
  \hspace{0.55cm}%
  \begin{minipage}{0.9\textwidth}
    \vspace{1em}
    \small
    \textbf{Evaluating Associations --- Real or Random?}\\[6pt]
    For each of the 3 scenarios on your activity sheet:
    \begin{enumerate}
      \setlength{\itemsep}{0.4em}\setlength{\topsep}{0em}
      \item Identify the two variables and their types.
      \item Decide: is this association likely real or possibly random?
            (Consider the sample size.)
      \item Name one confounding variable.
      \item Could this study establish causation?
    \end{enumerate}
    \vspace{0.6cm}
    {\color{redacc}\bfseries\small TIME: 8--10 minutes}
  \end{minipage}
\end{frame}
}

% ─────────────────────────────────────────────────────────────────────────────
% SLIDE 9 — EXIT TICKET
% ─────────────────────────────────────────────────────────────────────────────
{
\setbeamercolor{background canvas}{bg=navy}
\setbeamercolor{enumerate item}{fg=white}
\begin{frame}[t, plain]
  \begin{tikzpicture}[remember picture, overlay]
    \fill[navylight] (current page.north west)
      rectangle ([yshift=-1.35cm]current page.north east);
  \end{tikzpicture}
  \vspace{0.45cm}\par
  \hspace{0.55cm}%
  \begin{minipage}{0.9\textwidth}
    {\color{goldacc}\bfseries\footnotesize\MakeUppercase{Exit Ticket}}
    \vspace{0.4cm}\par
    {\color{skymid}\small\itshape
      A researcher surveys 30 high school students and finds that students
      who drink more coffee tend to have higher GPAs.}\par
    \vspace{0.4cm}
    \small
    \begin{enumerate}
      \setlength{\itemsep}{0.45em}\setlength{\topsep}{0em}
      \item {\color{white}What are the two variables? What type is each?}
      \item {\color{white}Could this association be random? Why does sample size matter?}
      \item {\color{white}Name one confounding variable.}
      \item {\color{white}Does this prove coffee causes higher grades? Explain.}
    \end{enumerate}
  \end{minipage}
\end{frame}
}

% ─────────────────────────────────────────────────────────────────────────────
% SLIDE 10 — CLOSING & UNIT 2 ROAD MAP
% ─────────────────────────────────────────────────────────────────────────────
{
\setbeamercolor{background canvas}{bg=navy}
\setbeamercolor{itemize item}{fg=white}
\begin{frame}[t, plain]
  \begin{tikzpicture}[remember picture, overlay]
    \fill[navylight] (current page.north west)
      rectangle ([yshift=-1.35cm]current page.north east);
  \end{tikzpicture}
  \vspace{0.45cm}\par
  \hspace{0.55cm}%
  \begin{minipage}{0.9\textwidth}
    {\color{goldacc}\bfseries\footnotesize\MakeUppercase{Closing \& Unit 2 Road Map}}
    \vspace{0.45cm}\par
    \begin{columns}[T]

      \begin{column}{0.54\textwidth}
        \small
        {\color{goldacc}\bfseries Key Reminders}\\[5pt]
        \begin{itemize}
          \setlength{\itemsep}{0.4em}\setlength{\topsep}{0em}
          \item {\color{white}\textbf{Apparent patterns may be random.}
            {\color{skymid} Small samples can look like big patterns by chance.}}
          \item {\color{white}\textbf{Association $\neq$ Causation.}
            {\color{skymid} A confounding variable may explain both.}}
          \item {\color{white}\textbf{The question drives the method.}
            {\color{skymid} Variable types determine which display and analysis to use.}}
        \end{itemize}
      \end{column}

      \begin{column}{0.42\textwidth}
        \small
        {\color{goldacc}\bfseries Where We Go Next}\\[5pt]
        \begin{itemize}
          \setlength{\itemsep}{0.35em}\setlength{\topsep}{0em}
          \item {\color{white}\textbf{2.2--2.3:} Two categorical variables
            $\to$ two-way tables, bar charts}
          \item {\color{white}\textbf{2.4--2.9:} Two quantitative variables
            $\to$ scatterplots, regression}
          \item {\color{white}\textbf{Unit 3:} Experimental design
            $\to$ establishing causation}
        \end{itemize}
        \vspace{0.5cm}
        {\color{skymid}\itshape\small Homework due next class.}
      \end{column}

    \end{columns}
  \end{minipage}
\end{frame}
}

\end{document}
